
{% Q3656730[D]

\needspace{8\baselineskip} 
\item \rtask \ponto{\pt} 
Assinale a alternativa que apresenta corretamente a estrutura de dados na Linguagem Python que garante uma coleção de elementos únicos, ou seja, sem duplicatas.  

\begin{answerlist}[label={\texttt{\Alph*}.},leftmargin=*]
    \ti \texttt{tuple}
    \ti \texttt{list}
    \ti \texttt{stack}
    \di \texttt{set}
    \ti \texttt{dict}
\end{answerlist}

    % Tema central: Estruturas de dados em Python, com ênfase em coleções que garantem unicidade de elementos (não permitem duplicatas).
    % 
    % Análise conceitual: Em provas para Analista de Banco de Dados, é importante dominar as diferenças fundamentais entre as principais estruturas de dados. O enunciado destaca: "coleção de elementos únicos, sem duplicatas". Esta é a pista essencial!
    % 
    % Resumo das estruturas:
    % 
    %     Tuple: Sequência ordenada e imutável. Permite elementos repetidos.
    %     List: Sequência ordenada e mutável. Aceita duplicatas.
    %     Stack: Segue princípio LIFO (Last In, First Out), geralmente implementada com listas, permitindo duplicatas.
    %     Set: Conjunto desordenado de elementos únicos. Não aceita duplicidade de elementos.
    % 
    % Justificativa da alternativa correta (D) Set:
    % 
    % O set foi desenvolvido especificamente para armazenar apenas elementos únicos. Se tentar inserir duplicatas, elas são ignoradas automaticamente.
    % 
    % Exemplo prático:
    % meu_conjunto = {1, 2, 3, 3}
    % 
    % O comando acima resultará em {1, 2, 3}, pois os conjuntos eliminam as duplicatas.
    % 
    % Análise das alternativas incorretas:
    % 
    %     A) Tuple: Apesar de imutável, permite duplicatas. Ex: (1, 1, 2).
    %     B) List: Permite e armazena vários elementos repetidos. Ex: [1, 2, 2, 3].
    %     C) Stack: Não existe como estrutura nativa (usa lista). Permite duplicatas e não garante unicidade.
    % 
    % Estrategicamente, preste atenção a palavras-chave como "elementos únicos" nas questões: Tuple e List costumam aparecer para confundir, pois são bem conhecidas, mas não atendem ao critério principal. Stack pode ser uma armadilha por ser menos comum, mas não se aplica.
    % 
    % Segundo autores como Mark Lutz (Python Programming, O’Reilly), o set é a estrutura correta quando se deseja unicidade. Assim, a alternativa D (Set) é a única opção correta nesta questão.
}




